\section{Evaluation Protocol}
\label{sec:protocol}

Each final policy is evaluated deterministically, without physics
randomisation and without early episode cuts. An episode \emph{succeeds}
if the robot never falls and ends at least 1\,m ($\times s$) from spawn
after 300 steps (6\,s). Every level is rolled out in two independent
passes of 4 episodes each, giving a per-level success rate per pass.

\paragraph{Uniform levels.} 2048 levels drawn uniformly from the level
space, the training distribution of DR. We report mean success with a 95\%
bootstrap confidence interval over levels.

\paragraph{Cross-fitted CVaR.} Robustness is the mean success on the worst
$\alpha$ fraction of levels ($\alpha\in\{0.1, 0.25\}$), in the spirit of
the adversarial evaluation of~\citet{rutherford2024noregrets}. Ranking and
scoring levels on the same episodes biases this estimate low through
selection noise, so we rank levels on one pass and score them on the other,
in both directions, and average. In a simulation with 2000 levels of true
success 0.5 and 4 episodes per level, same-pass \cvar{0.25} reads 0.19
while the cross-fitted estimate reads 0.49.

\paragraph{Level reliability.} The Pearson correlation of per-level success
between the two passes. A high value means difficulty is a property of the
level rather than of the episode, which is a precondition for any curation
method to have something to select. On our earlier control-point
heightfield space this correlation was only 0.02--0.17; on the parametric
pilot space it was 0.41--0.74 (Appendix~\ref{app:pilot}).

\paragraph{Difficulty sweep.} The same 2048 genomes shrunk toward flat by
$d\in\{0.25,0.5,0.75,1.0\}$. We report $d=0.5$ and $d=1.0$; $d=0.25$ is
solved by almost any policy trained near flat terrain and does not separate
methods. A flat-terrain reference ($d=0$) checks that a method has not
traded basic walking for terrain skills.
\todo{$d=0$ is not in \texttt{eval\_brax.py} yet; add it or drop this
sentence.}

\paragraph{Solvable subset.} Mean success on the uniform levels that at
least one evaluated policy solves, which removes levels no method can pass
from the denominator.
\todo{Computed offline from the per-level \texttt{.npz} outputs; define
``solves'' (any episode vs.\ $\geq50\%$) before reporting.}

\paragraph{Capability envelopes.} For each of seven features (noise,
slope up, slope down, stairs up, stairs down, gaps, blocks), the feature
ramps from off to the end of its range in six steps with all other
features off, 64 seeded levels per step. The \emph{envelope} is the
largest feature magnitude still passed with 50\% success (linearly
interpolated), the \emph{AUC} the mean success over the ramp.
Envelopes are in Wojtek units for both robots.

\paragraph{Held-out suite.} Seven structured terrain families no method
trains on, each at four strengths $k=1..4$ and 64 instances, radially
symmetric like the training levels (Wojtek units, $\times s$ for Go1):
stairs up and down ($2k$\,cm steps every 0.3\,m), slopes up and down
($5k^\circ$), trenches (10\,cm deep, $5k$\,cm wide, every 0.5\,m), waves
($\pm2k$\,cm, 0.6\,m wavelength) and random blocks ($\pm2k$\,cm, 0.3\,m
cells). We report mean success over the 28 cells and the CVaR over cells.

\paragraph{Sim2sim transfer.} Policies trained in MJX are re-evaluated in
CPU MuJoCo under five physics variants: \texttt{mjx\_match} (settings
matching training, a sanity check), \texttt{newton\_elliptic},
\texttt{half\_dt}, \texttt{full\_collision} and \texttt{soft\_contact}.
\todo{Describe each variant (solver, friction cone, timestep, collision
set, contact softness) from \texttt{src/scripts/eval\_sim2sim.py} once it
exists, and state which level sets it runs.}

\paragraph{Aggregation.} Per-seed metrics are averaged over three seeds;
per-level outputs are stored so that interquartile means and
stratified bootstrap intervals across seeds~\citep{agarwal2021precipice}
can be computed.
\todo{Decide between seed mean $\pm$ std and IQM with stratified bootstrap
for the main table; three seeds is few for either.}
